% =============================================================================
% cap2_hardware/03_acondicionamento_geometria.tex
% Seção 2.4 -- Geometria de captura e proteção (v2.2 — segunda compactação)
% =============================================================================
\section{Geometria de captura e proteção}
\label{sec:acondicionamento-geometria}

Qualquer que seja a tipologia escolhida na Seção~\ref{sec:tipologias-arquiteturais}, o módulo de aquisição deve sobreviver fisicamente ao ambiente do Campus 2 e produzir quadros cujas \emph{features} visuais sejam efetivamente utilizáveis pelas etapas seguintes. Esta seção trata das duas dimensões físicas transversais a essas tipologias --- acondicionamento dos componentes e geometria de captura junto aos abrigos --- no plano da análise de viabilidade.

A primeira preocupação de acondicionamento é \textbf{ambiental}. Os pontos P01--P10 estão expostos a chuva, orvalho, poeira, variação térmica diária acima de \(15\,^\circ\mathrm{C}\) e, em vários casos, contato direto com vegetação e solo úmido. Invólucros classe \textbf{IP66} (proteção total contra poeira e jatos d'água) ou \textbf{IP67} (imersão temporária) são requisito não negociável. Em paralelo, surge o \emph{trade-off} clássico entre vedação e ventilação: módulos com SBC e NPU dissipam calor não desprezível em operação contínua, e invólucros completamente vedados podem elevar a temperatura interna a ponto de comprometer a estabilidade dos sensores, problema reportado em sistemas semelhantes operados em climas tropicais~\cite{tiddeman2023lowpowerhardware,sargent2021raspberrypi}. Essa restrição impacta o regime térmico do sensor de imagem, com possíveis consequências sobre ruído, faixa dinâmica e estabilidade colorimétrica dos quadros entregues ao pipeline.

A segunda preocupação é \textbf{patrimonial e cooperativa}. \emph{Camera trapping} é, internacionalmente, campo afetado por furto e vandalismo: levantamentos com centenas de profissionais documentam perdas anuais expressivas e indicam que as estratégias mais efetivas combinam camuflagem visual, fixação robusta e posicionamento que reduza a saliência do equipamento~\cite{meek2018theft}. No Campus 2, o sistema será instalado em ambiente comunitário cooperativo, partilhado por estudantes, técnicos, voluntários e visitantes externos. Câmeras visivelmente identificáveis tendem a alterar tanto o comportamento de pessoas quanto, em alguma medida, a frequentação dos abrigos pelos próprios gatos --- componente comportamental relevante em estudos urbanos de fauna sinantrópica~\cite{cove2022capital}. A diretriz é, portanto, integrar o módulo às ``casinhas de reciclagem'' existentes, aproveitando seu volume para abrigar a caixa eletrônica. Para o pipeline, isso implica que a posição do sensor é dada pelo abrigo, convertendo-se em parâmetro fixo de calibração geométrica por ponto.

A geometria com que a câmera observa o ponto precede qualquer escolha de modelo, pois determina o que é geometricamente possível extrair do quadro. Para o pipeline aqui projetado, dois conjuntos de \emph{features} precisam ser preservados em fração estatisticamente significativa dos quadros: a \textbf{face} do animal (relevante para abordagens face-a-face habilitadas pelo \modeloPetFace) e o \textbf{flanco}, incluindo padrão de pelagem, contorno corporal e marcação de orelha (\emph{ear-tipping}), conforme o método de re-ID por partes do corpo de \citeonline{akbar2025bodypart}. Posicionamentos que registrem o animal apenas pelo dorso ou exclusivamente em silhueta contraluz comprometem a reidentificação antes mesmo da escolha do modelo, transferindo ao pipeline um problema geometricamente insolúvel.

A literatura fornece direção qualitativa robusta. Para mamíferos de pequeno e médio porte, à altura típica dos gatos comunitários, alturas baixas (algumas dezenas de centímetros, próximas ao nível dos ombros do animal) maximizam taxa de detecção e qualidade da identificação individual, ao passo que posicionamentos elevados, frequentemente adotados como medida antifurto, reduzem significativamente tanto a probabilidade de disparo quanto a utilidade dos quadros para classificação fina~\cite{seidlitz2020height,beukes2025refining,meek2016higher}. No plano angular, ângulos levemente picados, com linha de visada apontada para o entorno imediato dos potes, oferecem o melhor compromisso entre cobertura do cone útil, captura simultânea de face e flanco e minimização de obstrução por vegetação rasteira~\cite{beukes2025refining,moll2020viewshed}. A obstrução do cone de visão por vegetação no primeiro plano reduz substancialmente a taxa efetiva de detecção --- queda da ordem de duas a três vezes em sítios visualmente carregados~\cite{moll2020viewshed} ---, diretriz que se traduz, no Campus 2, em remoção periódica de folhagem incorporada à rotina dos cuidadores. A iluminação ativa por NIR precisa estar coalinhada com o eixo óptico e dimensionada para o alcance típico do ponto, sob pena de gerar quadros saturados próximos ao módulo e escuros para além de distância curta. Essas restrições geométricas e fotométricas são parâmetros de entrada do pipeline do Capítulo~\ref{cap:pipeline}: o que o software pode extrair sobre indivíduos da colônia está limitado, em última instância, pela qualidade angular, espacial e fotométrica dos quadros que a infraestrutura entrega.
